\chapter{Base SWEB vs.\ your team repo}
\label{ch:team}

The interview happens in \textbf{your group's repository}, not in base SWEB. Since Phase 0 it
separates process and thread. The mechanisms are identical; some names and the place for
per-process state differ. State of \code{osw26e3} \code{main} on 6 Oct 2026 --- check again
before the interview, the team keeps changing it.

\begin{center}
\small
\begin{tabular}{L{6.4cm}L{8.4cm}}
\toprule
\textbf{Base SWEB / this course} & \textbf{Team repo (osw26e3)} \\
\midrule
\code{class UserProcess : public Thread} (the process \emph{is} its only thread) &
\code{class UserProcess} (no thread), \code{class UserThread : public Thread}, \code{createMainThread()} \\
\code{currentThread->loader\_} & \code{currentThread->getLoader()} (the loader belongs to the process) \\
\code{currentThread->process\_} (added in A0-08) & \code{currentThread->getProcess()} (\code{Thread::process\_}) \\
per-process state in \code{Loader} (\code{heap\_size\_}, segv handler, \ldots) & per-process state in \code{UserProcess} \\
\code{Loader::arch\_memory\_lock\_} (added by this course) & whatever lock your team uses for the page tables --- \textbf{find it before the interview}; if there is none, that is the first thing a tutor asks about \\
\code{UserProcess::createThread} (A0-08) & your team's \code{pthread\_create} (P1) --- know its steps \\
syscall numbers 1400--1799 & your own block (1500--1999) --- avoid collisions with the team's syscalls \\
\code{process\_->waitForThreads()} in \code{exit} & your team's exit/join design \\
\bottomrule
\end{tabular}
\end{center}

\begin{infobox}[title=\textbf{Before the interview: 15 minutes in the team repo}]
\begin{enumerate}
\item Where is the page-table lock? Which locks does \code{UserProcess} have, what do they protect?
\item Where does \code{pthread\_create} map the stack, set the registers, add the thread?
\item Where are a process's threads listed, and who frees the address space when?
\item Add a trivial syscall in a scratch branch and run it --- so the six places are fresh.
\end{enumerate}
\end{infobox}
